
\chapter{Graph Neural Networks as EPMs}
\label{ch:gnns_epms}

\begin{figure}[t]
	\centering
	\includegraphics[width=0.75\textwidth]{figures/sat.pdf}
	\caption[Overview of SAT graph types]{Overview of the three graph types generated from a given \gls{sat} formula.
		\label{fig:sat}}
\end{figure}

The \new{Boolean satisfiability problem} (\gls{sat}) is a central and longstanding \textsf{NP}-complete problem~\citep{karp1972reducibility,biere_handbook_2021}. It asks whether there exists a truth assignment that satisfies a Boolean formula $\varphi$ over variables $v_1,\ldots,v_r$, typically expressed in \new{conjunctive normal form} (\gls{cnf}):
\begin{equation*}
	\varphi = \bigwedge_{i=1}^n c_i, \quad \text{where each } c_i = \bigvee_{j=1}^{m_i} l_{i, j}.
\end{equation*}
Here, each $l_{i,j}$ is a literal, i.e., either a variable $v_p$ or its negation $\lnot v_p$. The task is to decide whether there exists an assignment $a$ such that $\varphi(a) = \text{true}$. \Gls{sat} is of significant theoretical importance as one of the first problems proven \textsf{NP}-complete~\citep{Cook71}, and it underpins numerous practical applications in hardware and software verification~\citep{BiereEtAl09}, automated planning~\citep{KautzSelman96}, and operations research~\citep{GomesEtAl08}. Its significance has led to the development of a wide range of solvers, systematically benchmarked in annual \gls{sat} competitions (e.g.,~\citep{FroEtAl21,BalEtAl17,HeuEtAl19}). A key property of modern \gls{sat} solvers is \textit{performance complementarity}---no single solver dominates across all instances. This motivates the study of \emph{algorithm selection}~\citep{Rice76,XuEtAl08}, where the aim is to choose the most effective solver for a given instance. \new{Empirical performance prediction} is a closely related problem, which seeks to forecast the computation time (or another performance metric) of a solver on a given instance. Such predictions can support algorithm selection, configuration, scheduling, and explainability~\citep{HutterEtAl14}. State-of-the-art methods for algorithm selection and performance prediction rely on \textit{hand-crafted features} extracted from \gls{sat} instances. The most widely adopted feature set is that of the \textsc{SATzilla} family of algorithm selectors~\citep{NudEtAl04,HutterEtAl14,ShaHoo24}. These features include basic structural properties (e.g., number of clauses and variables), graph-based descriptors derived from instance structure, and probing features obtained by briefly running a solver to collect computation-time statistics. Machine learning models---most often tabular models---are then trained on these features for prediction or selection tasks.

\Gls{sat} instances can be naturally represented as graphs, reflecting their permutation-invariant structure that arises from the commutativity and associativity of logical conjunction and disjunction. Indeed, many of the hand-crafted \textsc{SATzilla} features are derived from such graph representations, including the
\begin{itemize}
	\item \new{Variable-clause graph}: a bipartite, undirected graph with a node for each variable $v$ and each clause $c$, where an edge connects a variable to a clause if and only if the variable appears in that clause.

	\item \new{Variable graph}: an undirected graph with one node per variable, where two variables $v_i$ and $v_j$ are connected if they co-occur in at least one clause.

	\item \new{Clause graph}: an undirected graph with one node per clause, where two clauses $c_i$ and $c_j$ are connected if they share at least one negated literal.
\end{itemize}

\textsc{SATzilla} extracts statistical properties of these graphs---such as mean degree, coefficient of variation, clustering coefficient, and graph diameter---and uses them as features for downstream learning tasks.

Several studies have explored the use of \glspl{gnn} or MPNNs to predict the satisfiability of a formula, e.g.,~\citet{SelEtAl19}. Some of this work employs an additional representation, the \new{literal-clause graph}, a bipartite graph with one node per literal (a variable or its negation) and one node per clause, with edges linking literals to the clauses they appear in. However, many such approaches are \emph{unsound}---they can produce incorrect results---or cannot guarantee valid satisfying assignments or proofs of unsatisfiability, limiting their practical usefulness~\citep{SelEtAl19,LiEtAl24}.  Other research has integrated MPNNs into existing \gls{sat} solvers to replace or guide solver heuristics~\citep{WanEtAl24,TonGro25}. These approaches retain the soundness and completeness guarantees of classical solvers while potentially accelerating the solving process. Nonetheless, benchmarking novel MPNN architectures in this setting is challenging because it requires extensive pretraining and large-scale evaluation.  A complementary line of work applies MPNNs to algorithm selection in \gls{sat}~\citep{ZhaEtAl24,Shavit23}, predicting the most effective solver for a given instance directly from graph-based representations.

We add node features to provide the \gls{gnn} with additional information about the formula's structure. Our node embeddings are inspired by the hand-crafted features used in \textsc{SATzilla}~\citep{NudEtAl04} and the node-level encodings proposed in \textsc{GraSS}~\citep{ZhaEtAl24}. These features capture local syntactic properties (e.g., polarity counts, clause membership statistics), normalized and relational quantities (e.g., positive-to-negative ratio, normalized frequency), and sinusoidal positional encodings for clause nodes. \Cref{tab:vcg_features,tab:lcg_features,tab:vg_features} summarize the node features for each graph type. All features are computed during graph construction and stored as floating-point node attributes.

\begin{table}[htbp]
	\centering
	\caption[VCG node features]{Node features for the variable-clause graph (VCG). Variable nodes carry 12 features, and clause nodes carry 22 features. Edges between variables and clauses are attributed with the literal polarity ($+1$/$-1$); reverse edges are included.}
	\label{tab:vcg_features}
	\renewcommand{\arraystretch}{1.2}
	\resizebox{\columnwidth}{!}{%
		\begin{tabular}{@{} p{0.22\textwidth} c p{0.62\textwidth} @{}}
			\toprule
			\textbf{Node type} & \textbf{Dim.} & \textbf{Feature description}                                              \\
			\midrule
			\multirow{11}{=}{\textbf{Variable}\newline(12 features)}
			                   & 1--2          & Constant placeholders (reserved, set to zero)                             \\
			                   & 3             & Number of Horn clauses in which this variable appears                     \\
			                   & 4             & Number of occurrences as a positive literal                               \\
			                   & 5             & Number of occurrences as a negative literal                               \\
			                   & 6             & Total number of clause appearances (degree)                               \\
			                   & 7             & Ratio of positive to negative occurrences                                 \\
			                   & 8             & Number of appearances in binary clauses (length two)                      \\
			                   & 9             & Number of appearances in unit clauses (length one)                        \\
			                   & 10            & Number of appearances in ternary clauses (length three)                   \\
			                   & 11            & Normalised frequency (degree divided by the total number of clauses)      \\
			                   & 12            & Mean length of clauses containing this variable                           \\
			\midrule
			\multirow{13}{=}{\textbf{Clause}\newline(22 features)}
			                   & 1--10         & Sinusoidal positional encoding (10 dimensions)                            \\
			                   & 11            & Number of positive literals                                               \\
			                   & 12            & Number of negative literals                                               \\
			                   & 13            & Ratio of positive to negative literals                                    \\
			                   & 14            & Clause length (number of literals)                                        \\
			                   & 15            & Indicator: unit clause (length one)                                       \\
			                   & 16            & Indicator: binary clause (length two)                                     \\
			                   & 17            & Indicator: ternary clause (length three)                                  \\
			                   & 18            & Indicator: Horn clause (at most one positive literal)                     \\
			                   & 19            & Normalized clause length (divided by the maximum clause length)           \\
			                   & 20            & Normalized clause position (index divided by the total number of clauses) \\
			                   & 21            & Indicator: definite Horn clause (exactly one positive literal)            \\
			                   & 22            & Indicator: negative clause (all literals are negative)                    \\
			\bottomrule
		\end{tabular}}
\end{table}

\begin{table}[htbp]
	\centering
	\caption[LCG node features]{Node features for the literal-clause graph (LCG). Literal nodes carry 12 features, and clause nodes carry 22 features. Literal-to-clause edges carry a polarity attribute ($+1$/$-1$); reverse and complement edges (linking each literal to its negation) are also included.}
	\label{tab:lcg_features}
	\renewcommand{\arraystretch}{1.2}
	\resizebox{\columnwidth}{!}{%
		\begin{tabular}{@{} p{0.22\textwidth} c p{0.62\textwidth} @{}}
			\toprule
			\textbf{Node type} & \textbf{Dim.} & \textbf{Feature description}                                              \\
			\midrule
			\multirow{12}{=}{\textbf{Literal}\newline(12 features)}
			                   & 1             & Indicator: positive literal                                               \\
			                   & 2             & Indicator: negative literal                                               \\
			                   & 3             & Number of Horn clauses in which this literal appears                      \\
			                   & 4             & Total number of clause occurrences                                        \\
			                   & 5             & Number of clause occurrences of the complementary literal                 \\
			                   & 6             & Degree (equal to the occurrence count)                                    \\
			                   & 7             & Ratio of this literal's occurrences to those of its complement            \\
			                   & 8             & Number of appearances in binary clauses (length two)                      \\
			                   & 9             & Number of appearances in unit clauses (length one)                        \\
			                   & 10            & Number of appearances in ternary clauses (length three)                   \\
			                   & 11            & Normalized frequency (occurrences divided by the total number of clauses) \\
			                   & 12            & Mean length of clauses containing this literal                            \\
			\midrule
			\multirow{13}{=}{\textbf{Clause}\newline(22 features)}
			                   & 1--10         & Sinusoidal positional encoding (10 dimensions)                            \\
			                   & 11            & Number of positive literals                                               \\
			                   & 12            & Number of negative literals                                               \\
			                   & 13            & Ratio of positive to negative literals                                    \\
			                   & 14            & Clause length (number of literals)                                        \\
			                   & 15            & Indicator: unit clause (length one)                                       \\
			                   & 16            & Indicator: binary clause (length two)                                     \\
			                   & 17            & Indicator: ternary clause (length three)                                  \\
			                   & 18            & Indicator: Horn clause (at most one positive literal)                     \\
			                   & 19            & Normalized clause length (divided by the maximum clause length)           \\
			                   & 20            & Normalized clause position (index divided by the total number of clauses) \\
			                   & 21            & Indicator: definite Horn clause (exactly one positive literal)            \\
			                   & 22            & Indicator: negative clause (all literals are negative)                    \\
			\bottomrule
		\end{tabular}}
\end{table}

\begin{table}[htbp]
	\centering
	\caption[VG node features]{Node features for the variable graph (VG). Each variable node carries 12 features. Edges connect co-occurring variables and carry no attributes.}
	\label{tab:vg_features}
	\renewcommand{\arraystretch}{1.2}
	\resizebox{\columnwidth}{!}{%
		\begin{tabular}{@{} p{0.22\textwidth} c p{0.62\textwidth} @{}}
			\toprule
			\textbf{Node type} & \textbf{Dim.} & \textbf{Feature description}                                                                \\
			\midrule
			\multirow{12}{=}{\textbf{Variable}\newline(12 features)}
			                   & 1             & Number of Horn clauses in which this variable appears                                       \\
			                   & 2             & Number of occurrences as a positive literal                                                 \\
			                   & 3             & Number of occurrences as a negative literal                                                 \\
			                   & 4             & Total co-occurrence count (degree in the variable graph)                                    \\
			                   & 5             & Ratio of positive to negative occurrences                                                   \\
			                   & 6             & Number of appearances in binary clauses (length two)                                        \\
			                   & 7             & Number of appearances in unit clauses (length one)                                          \\
			                   & 8             & Number of appearances in ternary clauses (length three)                                     \\
			                   & 9             & Normalized frequency (clause appearances divided by the total number of clauses)            \\
			                   & 10            & Mean length of clauses containing this variable                                             \\
			                   & 11            & Total number of distinct clauses containing this variable                                   \\
			                   & 12            & Mean co-occurrence degree (co-occurrence count divided by the number of containing clauses) \\
			\bottomrule
		\end{tabular}}
\end{table}

In \gb, we introduce, to the best of our knowledge, the largest algorithm-selection benchmark for \gls{sat} solvers, covering eleven solvers and more than 100\,000 instances. Our benchmark supports multiple graph representations of \gls{sat} formulae and also provides the \textsc{SATzilla} 2024 feature set~\citep{ShaHoo24}, enabling their combination with MPNNs. Unlike existing \gls{gnn} benchmarks for \gls{sat}~\citep{LiEtAl24}, which primarily focus on satisfiability prediction or assignment generation, our benchmark is designed for practical downstream tasks in \gls{sat} solving---specifically, algorithm selection and performance prediction.
\cref{fig:sat} provides an overview of how the three graph types are constructed from a given Boolean formula.


\paragraph{Related work}~
As the use of MPNNs for algorithm selection is a relatively new research direction, no benchmarks exist for this task. Previously, \citet{Shavit23} used a dataset of 3\,000 synthetically generated \gls{sat} instances, while \citet{ZhaEtAl24} used a proprietary dataset as well as the 2022 Anniversary track of the \gls{sat} Competition, while removing large instances. Outside of the graph domain, the \textsc{ASlib} benchmark suite is used to benchmark algorithm selection methods~\citep{BisEtAL16}. \textsc{ASlib} contains multiple algorithm-selection datasets across various domains, including \gls{sat} and TSP. Each dataset contains features extracted from instances (such as \textsc{SATzilla}) as well as the computation times of several algorithms on those instances. We note that \textsc{ASlib} provides \gls{sat} scenarios based on a subset of \gls{sat} competitions, including instances and solvers from specific years. In contrast, our dataset comprises all available \gls{sat} competition instances, as well as numerous instances from other sources.
% }

\paragraph{Description of the learning task}~
We introduce two learning tasks. \new{Performance prediction} is a regression problem that aims to predict the computation time of \gls{sat} solvers on unseen instances, arising in applications such as algorithm selection \citep{Rice76} and algorithm configuration \citep[e.g.][]{LinEtAl22}. \new{Algorithm selection} is a multi-class classification problem that aims to select the best algorithm for a given instance.
While algorithm selection can be reduced to performance prediction by choosing the solver with the lowest predicted computation time, we treat them as distinct problems. For algorithm selection, we adopt the loss function proposed by~\citep{ZhaEtAl24}, which penalizes the predicted probabilities in proportion to the solver's computation time. Intuitively, a solver that only slightly underperforms on an instance incurs a small penalty, while poor selections are penalized more heavily, i.e.,
\begin{equation*}
	\frac{1}{N} \sum_{i=1}^N \left( \sum_{k=1}^K p_i^k \cdot t_i^k - t_i^* \right)^2,
\end{equation*}
where $p_i^k$ is the predicted probability of selecting solver $k$ for instance $i$, $t_i^k$ is the computation time of solver $k$ on instance $i$, and $t_i^* = \min_k t_i^k$ is the computation time of the \new{virtual best solver} (\gls{vbs}). Each instance is provided with
three graph-based representations VG, VCG, and LCG, described above.
In all cases, \textsc{SATzilla} features can be integrated as node attributes to enrich the graph representation.

The performance of a solver $s$ on an instance set $\Pi=\{\pi_1, \pi_2, \dots, \pi_n\}$ is measured using the widely adopted \new{penalized average computation time} ($\text{\gls{par}}_k$) metric, where $k$ denotes the penalty factor for unsolved instances under a cutoff time $c$. The $\text{\gls{par}}_k(s, \Pi)$ score is defined as the mean computation time of $s$ across $\Pi$, where each unsolved instance contributes a computation time of $k \cdot c$. Lower $\text{\gls{par}}_k$ values indicate better performance. For algorithm selectors, both feature extraction and model inference time are included in the total solving time, implying that MPNN inference must be efficient to be practically helpful.

For the performance prediction task, the model output is the base-10 logarithm of the $\text{\gls{par}}_{10}$ score of solver $s$,
\begin{equation*}
	\log_{10}(\text{\gls{par}}_{10}(s, \Pi)),
\end{equation*}
a transformation shown to improve predictive accuracy in prior work~\citep{HutterEtAl14}. Model performance is evaluated using the \new{root mean-squared error} (\gls{rmse}) between predicted and actual log-scaled computation times.

For algorithm selection, we use the \new{closed gap} (\gls{cg}) metric, which quantifies how close a selector comes to the performance of the \gls{vbs} relative to the single best solver (\gls{sbs}). Formally,
\begin{equation*}
	\text{\gls{cg}}(\Pi) = \frac{\text{\gls{par}}_{10}(\text{\gls{sbs}}, \Pi) - \text{\gls{par}}_{10}(s, \Pi)}{\text{\gls{par}}_{10}(\text{\gls{sbs}}, \Pi) - \text{\gls{par}}_{10}(\text{\gls{vbs}}, \Pi)}.
\end{equation*}
A higher \gls{cg} score indicates stronger selection performance, with \gls{cg} $=1$ corresponding to \gls{vbs} performance and \gls{cg} $=0$ to \gls{sbs} performance.

\paragraph{Details on the datasets}~
We created the largest algorithm-selection scenario for \gls{sat} solvers to date. We collected all \gls{sat} Competition instances via GBD~\citep{IseJab24} and added all instances from the AClib benchmark for algorithm configuration~\citep{HutEtAl14b}. Together, these sources include \gls{sat} formulae from diverse real-world applications and synthetically generated instances. To further enrich the dataset, we generated additional formulae using three well-established \gls{sat} instance generators, i.e.,

\begin{itemize}
	\item \new{FuzzSAT}~\citep{BruEtAl10} produces \gls{cnf} fuzzing instances originally designed to identify issues in digital circuits. It generates random Boolean circuits as DAGs and converts them into \gls{cnf} via the Tseitin transformation~\citep{tseitin1983complexity}. To construct the \gls{dag}, we used $v \in [50,150]$ input variables; operands were randomly paired and connected by randomly chosen operator nodes until each input was referenced at least $t \in [1,5]$ times. Finally, random clauses of length $s \in [3,8]$ were added to increase complexity.

	\item \new{Cryptography}~\citep{SaeEtAl17,AlaEtAl24} generates \gls{sat} instances encoding preimage attacks on MD4, SHA-1, and SHA-256. Each encoding applies multiple \emph{rounds} of operations; higher numbers of rounds yield harder instances. We set the number of rounds to $d \in [16,30]$, which yields formulae that are challenging for modern solvers but not unsolvable. Additionally, we fixed $i \in [0,384]$ of the 512 input bits; fixing more bits results in easier instances.

	\item \new{Community attachment}~\citep{GirLev16} produces pseudo-random \gls{sat} formulae with community structure, a property often observed in real-world \gls{sat} encodings (e.g., hardware and software verification tasks~\citep{AnsEtAl12}). The generator can create relatively small but difficult instances. We used $n \in [1\,000, 3\,000]$ variables, $c \in [5,100]$ communities, a clause-to-variable ratio of $[3{.}7, 4{.}5]$, and modularity factor $q \in [0{.}3,0{.}9]$. Higher modularity values enforce stronger community structure.
\end{itemize}

For each generator, parameter ranges were chosen to yield challenging but still solvable instances.  To ensure this, we excluded ranges of values yielding no instances solvable by \solver{Kissat}, the winner of the 2024 \gls{sat} Competition, between 1 and 5\,000 seconds. After generation, we applied the \textsc{Satellite} preprocessor to remove redundant clauses and variables. A summary of all generated instances is provided in \cref{tab:sat_generatred_instances}.

We pre-selected solvers for the benchmark by constructing a portfolio of $m$ complementary state-of-the-art \gls{sat} solvers, based on the results of the 2024 \gls{sat} Competition. Portfolios were built using beam search to identify the set of $m$ solvers achieving the \gls{vbs} performance. To determine $m$, we evaluated portfolios of all possible sizes and plotted portfolio size against achieved \gls{vbs}. We then identified the inflection point (the ``knee'') of this curve using the Kneedle algorithm~\citep{SatEtlA11} and adopted the corresponding portfolio. The \gls{asf} library~\url{https://github.com/hadarshavit/asf} was used to perform the selection procedure. The same procedure was repeated on the \gls{sat} Competition 2023.

In total, our dataset includes the following solvers. From the 2024 \gls{sat} Competition: Kissat~\citep{BieEtAl24}, BreadIdKissat~\citep{BogEtAl24}, AMSAT~\citep{LiEtAl24b}, and KissatMABDC~\citep{LiuEtAl24}. From the 2023 \gls{sat} Competition: SBVA CadiCal~\citep{HabEtAl23}, Reencode Kissat~\citep{ReeEtAl23}, Minisat XOR, KissatMAB Prop PR~\citep{GaoEtAl23}, hKissatInc~\citep{KonEtAl23}, CadiCal vivinst~\citep{BieEtAl23}, and BreakID Kissat~\citep{BogEtAl23}. Solver computation times were measured in CPU time using \texttt{runsolver}~\citep{Roussel11} and the generic wrapper tools~\citep{EggEtAl19}, ensuring accurate and consistent performance measurements.


\begin{table}
	\centering
	\caption[SAT solving benchmark instances]{Overview of the instances for the \gls{sat} solving benchmark.}
	\resizebox{\columnwidth}{!}{
		\begin{tabular}{ll@{\hskip 0.5in}llll@{\hskip 0.5in}llll}
			\toprule
			\textbf{Source}      & \textbf{\# of Instances} & \multicolumn{4}{c}{\textbf{\# of Variables}} & \multicolumn{4}{c}{\textbf{\# of Clauses}}                                                                                   \\
			                     &                          & min                                          & max                                        & avg        & median     & min    & max           & avg            & median      \\
			\midrule
			\Gls{sat} Competition      & 31\,024                  & 3                                            & 25\,870\,369                               & 69\,190.77 & 13\,168.50 & 7      & 871\,935\,536 & 1\,022\,025.50 & 118\,526.00 \\
			Community Attachment & 29\,994                  & 922                                          & 2\,956                                     & 1\,931.05  & 1\,928.00  & 3\,566 & 13\,413       & 8\,088.71      & 8\,041.00   \\
			AClib                & 33\,955                  & 6                                            & 248\,738                                   & 3\,319.65  & 1\,270.00  & 24     & 103\,670\,100 & 136\,471.03    & 10\,000.00  \\
			FuzzSAT              & 8\,020                   & 2                                            & 944\,685                                   & 10\,217.75 & 1\,820.00  & 1      & 4\,094\,516   & 45\,638.30     & 8\,074.50   \\
			Cryptography         & 4\,873                   & 6                                            & 12\,415                                    & 4\,264.89  & 3\,225.00  & 24     & 185\,006      & 55\,309.63     & 45\,316.00  \\\midrule
			\textbf{Total}       & 107\,866                 & 2                                            & 25\,870\,369                               & 22\,434.71 & 1\,879.00  & 1      & 871\,935\,536 & 345\,051.72    & 10\,177.50  \\
			\bottomrule
		\end{tabular}
	}
	\label{tab:sat_generatred_instances}
\end{table}
We provide three datasets, \dataset{small}, which includes only small formulae with up to 3\,000 variables and 15\,000 clauses, \dataset{medium}, which consists of the small formulae with additional medium size formulae with size of up to 20\,000 variables and 80\,000 clauses and \dataset{large}, which includes all formulae. This way, we provide datasets with varying levels of hardness, as formula size can pose hurdles for training time and GPU memory. While the \dataset{large} and \dataset{medium} datasets are unlikely to be accessible on current hardware, we provide them to keep the benchmark relevant in the long term.


\paragraph{Dataset challenges:} Algorithm selection and performance prediction are particularly challenging in this setting for several reasons. First, solver performance is inherently stochastic: repeated runs of the same solver on the same instance can result in substantially different running times. This induces noisy target values that the learning algorithm must account for. Second, the targets are right-censored due to the use of solver cutoffs: when a solver does not terminate before the cutoff, its true running time is only known to exceed the cutoff. Third, the \gls{sat} instances can be very large, with corresponding graphs containing thousands to millions of variable nodes, making scalable graph processing a major challenge. Finally, accurate prediction requires the \gls{gnn} to capture both local structural patterns and the formula's global properties.

\paragraph{Experimental evaluation and baselines}~
All \gls{sat} solving tasks in the benchmark are graph-level. For the graph transformer baselines, we employ an encoder-processor-decoder pipeline, as described in~\Cref{Section:ExpSetup}. We used the following architectures: GIN, GCN, GIN+, GCN+, Gated GCN+. While we considered GT for the tasks, we encountered high computational costs in calculating the positional embeddings (more than a CPU day) and therefore excluded them from the evaluation. To ensure a fair comparison with hand-crafted features and to evaluate the capacity of \glspl{gnn} to extract relevant information purely from the formula's topology, we did not provide the models with any initial node-level features, forcing them to learn solely from the graph structure. Furthermore, when training on medium- and large-sized graphs, we encountered out-of-memory errors even on an H100. Consequently, we only consider small formulae for MPNNs. Similarly, we observed out-of-memory issues with the clause graphs and therefore excluded them from the evaluation. We optimized the hyperparameters of these models using SMAC3~\citep{LinEtAl22} with a multi-fidelity facade, running 60 trials. Full details of the hyperparameter configuration space and the tuning procedure can be found in~\Cref{sec:hpo}.

We note that traditionally, algorithm selection datasets are evaluated using cross-validation, primarily due to limited data availability. In contrast, our new dataset contains tens of thousands of formulas. Since it is also intended for deep learning approaches, which are computationally expensive and thus impractical to evaluate using 10-fold cross-validation, we instead provide a fixed train–validation–test split and report results across 3 random seeds.

For the performance prediction task, we compare the MPNNs against traditional baselines, which use the \solver{SATzilla} features and a (tabular) machine learning model. We use random forest~\citep{Breiman01} and XGBoost~\citep{CheEtAl16}, which are commonly used as empirical performance models (\gls{epm}) (see, e.g., \citet{BanEtAl22,EggEtAl15,HutEtAl14}).
We use the \solver{SATzilla} 2024 features as input and predict the log10-scaled computation time. These baselines are implemented using the \gls{asf} library~\url{https://github.com/hadarshavit/asf}.
For algorithm selection, we provide several well-known baselines. First, we use pairwise regression~\citep{KerEtAl18} (PR; predicts performance differences between each pair of solvers, then picks the solver with the best sum), pairwise classification~\citep{XuEtAl12} (PC; predicts the better solver in each pair, then uses majority vote, %\mia{
winner of the ICON challenge on algorithm selection~\citep{LinEtAl19}%}
) as well as a simple multi-class classification~\citep{KotEtAl13} (MC; directly predicts the best solver). Similar to the \gls{epm} baselines, these baselines use the \solver{SATzilla} 2024 features with a random forest. We tune these baselines using SMAC3~\citep{LinEtAl22} for 50 trials.


\begin{table}[htbp]
	\tabledefaults
	\centering
	\caption[SAT solving RMSE on small instances]{\textbf{\gls{sat} Solving.} \gls{rmse} (log\textsubscript{10} time) for performance prediction on \dataset{small} \gls{sat} instances. We report results on three different graph representations (\dataset{VG}, \dataset{VCG}, \dataset{LCG}) and using the hand-crafted SATzilla features. Lower is better.}
	\label{table:sat_epm_res_small}
	\resizebox{0.6\textwidth}{!}{%
		\begin{tabular}{llcccc}
			\toprule
			\textbf{Solver} & \textbf{Method} & \dataset{VG}         & \dataset{VCG}        & \dataset{LCG}        & \dataset{SATzilla}   \\
			\midrule
			\multirow{8}{*}{\solver{Kissat}}
			                & GIN             & \meanstd{0.66}{0.02} & \meanstd{0.64}{0.07} & \meanstd{0.70}{0.01}                        \\
			                & GCN             & \meanstd{0.98}{0.00} & \meanstd{0.70}{0.01} & \meanstd{0.82}{0.07}                        \\
			                & GCN+            & \meanstd{0.69}{0.03} & \meanstd{0.51}{0.02} & \meanstd{0.67}{0.09}                        \\
			                & GIN+            & \meanstd{0.78}{0.02} & \meanstd{0.87}{0.05} & \meanstd{0.83}{0.07}                        \\
			                & Gated GCN+      & --                   & \meanstd{1.05}{0.04} & \meanstd{1.04}{0.04}                        \\
			                & RF              &                      &                      &                      & \meanstd{0.65}{0.00} \\
			                & XGB             &                      &                      &                      & \meanstd{0.64}{0.00} \\
			\midrule
			\multirow{8}{*}{\solver{BreakIDKissat}}
			                & GIN             & \meanstd{0.65}{0.01} & \meanstd{0.63}{0.04} & \meanstd{0.68}{0.01}                        \\
			                & GCN             & \meanstd{0.96}{0.00} & \meanstd{0.70}{0.01} & \meanstd{0.80}{0.09}                        \\
			                & GCN+            & \meanstd{0.69}{0.01} & \meanstd{0.58}{0.03} & \meanstd{0.62}{0.03}                        \\
			                & GIN+            & \meanstd{0.82}{0.10} & \meanstd{0.89}{0.03} & \meanstd{0.95}{0.08}                        \\
			                & Gated GCN+      & --                   & \meanstd{0.94}{0.02} & \meanstd{0.94}{0.03}                        \\
			                & RF              &                      &                      &                      & \meanstd{0.68}{0.00} \\
			                & XGB             &                      &                      &                      & \meanstd{0.67}{0.00} \\
			\midrule
			\multirow{8}{*}{\solver{KissatMABDC}}
			                & GIN             & \meanstd{0.64}{0.03} & \meanstd{0.72}{0.01} & \meanstd{0.70}{0.02}                        \\
			                & GCN             & \meanstd{0.99}{0.00} & \meanstd{0.75}{0.02} & \meanstd{0.74}{0.01}                        \\
			                & GCN+            & \meanstd{0.74}{0.07} & \meanstd{0.86}{0.01} & \meanstd{0.78}{0.07}                        \\
			                & GIN+            & \meanstd{0.83}{0.09} & \meanstd{1.05}{0.06} & \meanstd{0.94}{0.01}                        \\
			                & Gated GCN+      & --                   & \meanstd{0.91}{0.15} & \meanstd{1.05}{0.18}                        \\
			                & RF              &                      &                      &                      & \meanstd{0.70}{0.00} \\
			                & XGB             &                      &                      &                      & \meanstd{0.69}{0.00} \\
			\midrule
			\multirow{8}{*}{\solver{AMSAT}}
			                & GIN             & \meanstd{0.71}{0.02} & \meanstd{0.70}{0.01} & \meanstd{0.71}{0.01}                        \\
			                & GCN             & \meanstd{1.01}{0.00} & \meanstd{0.82}{0.10} & \meanstd{0.75}{0.01}                        \\
			                & GCN+            & \meanstd{0.85}{0.09} & \meanstd{0.61}{0.02} & \meanstd{0.74}{0.06}                        \\
			                & GIN+            & \meanstd{0.92}{0.07} & \meanstd{1.05}{0.04} & \meanstd{1.04}{0.06}                        \\
			                & Gated GCN+      & --                   & \meanstd{1.02}{0.01} & \meanstd{1.03}{0.07}                        \\
			                & RF              &                      &                      &                      & \meanstd{0.74}{0.00} \\
			                & XGB             &                      &                      &                      & \meanstd{0.74}{0.00} \\
			\bottomrule
		\end{tabular}}
\end{table}

\begin{table}[htbp]
	\tabledefaults
	\centering
	\caption[SAT solving RMSE on medium and large instances]{\textbf{\gls{sat} Solving.} \gls{rmse} (log\textsubscript{10} time) for performance prediction on \dataset{medium} and \dataset{large} \gls{sat} instances. Lower is better.}
	\label{table:sat_epm_res_med_large}
	\begin{minipage}{0.48\textwidth}
		\centering
		\subcaption{Medium-size \gls{sat} instances}
		\resizebox{0.8\linewidth}{!}{%
			\begin{tabular}{llc}
				\toprule
				\textbf{Solver} & \textbf{Method} & \dataset{SATzilla}   \\
				\midrule
				\multirow{2}{*}{\solver{Kissat}}
				                & RF              & \meanstd{0.63}{0.00} \\
				                & XGB             & \meanstd{0.62}{0.00} \\
				\midrule
				\multirow{2}{*}{\solver{BreakIDKissat}}
				                & RF              & \meanstd{0.66}{0.00} \\
				                & XGB             & \meanstd{0.65}{0.00} \\
				\midrule
				\multirow{2}{*}{\solver{KissatMABDC}}
				                & RF              & \meanstd{0.68}{0.00} \\
				                & XGB             & \meanstd{0.67}{0.00} \\
				\midrule
				\multirow{2}{*}{\solver{AMSAT}}
				                & RF              & \meanstd{0.70}{0.00} \\
				                & XGB             & \meanstd{0.69}{0.00} \\
				\bottomrule
			\end{tabular}}
		\label{table:sat_epm_res_med}
	\end{minipage}
	\hfill
	\begin{minipage}{0.48\textwidth}
		\centering
		\subcaption{Large \gls{sat} instances}
		\resizebox{0.8\linewidth}{!}{%
			\begin{tabular}{llc}
				\toprule
				\textbf{Solver} & \textbf{Method} & \dataset{SATzilla}   \\
				\midrule
				\multirow{2}{*}{\solver{Kissat}}
				                & RF              & \meanstd{0.64}{0.00} \\
				                & XGB             & \meanstd{0.64}{0.00} \\
				\midrule
				\multirow{2}{*}{\solver{BreakIDKissat}}
				                & RF              & \meanstd{0.66}{0.00} \\
				                & XGB             & \meanstd{0.65}{0.00} \\
				\midrule
				\multirow{2}{*}{\solver{KissatMABDC}}
				                & RF              & \meanstd{0.67}{0.00} \\
				                & XGB             & \meanstd{0.66}{0.00} \\
				\midrule
				\multirow{2}{*}{\solver{AMSAT}}
				                & RF              & \meanstd{0.69}{0.00} \\
				                & XGB             & \meanstd{0.69}{0.00} \\
				\bottomrule
			\end{tabular}}
		\label{table:sat_epm_res_large}
	\end{minipage}
\end{table}

\paragraph{Results}~
We start by presenting the results for the performance prediction task. We present the \gls{rmse} of the log-scaled computation time for the solvers from \gls{sat} Competition 2024 (\solver{Kissat}, \solver{BreakIDKissat}, \solver{KissatMABDC}, \solver{AMSAT}) for the small formulae in table~\Cref{table:sat_epm_res_small}. Note that missing values (--) for Gated GCN Plus with the VG representation indicate that the model encountered out-of-memory (OOM) errors on the GPU.
The table shows that the best performance for \solver{Kissat}, \solver{BreakIDKissat}, and \solver{AMSAT} is achieved by GCN Plus with the VCG graph, as this model outperforms the SATzilla features. For \solver{KissatMABDC}, a simple GIN with a variable graph performs best. We also note that the \gls{sat}-specific architecture, GraSS, performs surprisingly poorly across the board compared to general-purpose architectures.
To the best of our knowledge, this is the first time the SATzilla features have been outperformed on the performance prediction task. We furthermore present the results of random forest and XGBoost on medium-size and large formulae in~\Cref{table:sat_epm_res_med,table:sat_epm_res_large}, such that they can be compared against graph machine learning solutions in the future. Overall, we observe similar performance between XGBoost and random forest, with a slight advantage to XGBoost.

For the algorithm selection task, we report the
closed gap metric of different selectors in~\Cref{table:sat_as_res}.
For the small instances, we observe that gated GCN plus performs best, followed by GCN plus and GIN, all of which are trained on literal-clause graphs. These graph-based approaches outperform the SATzilla feature-based baselines, with only pairwise regression achieving a positive closed gap. Comparing the different graph representations across tasks, we observe clear trade-offs: while the VCG representation yields the strongest performance for the prediction task, the LCG representation enables the highest gap closure for algorithm selection. This suggests that the explicit literal polarity and complementary edges in LCG might be more beneficial for relative performance ranking than for absolute runtime prediction.
Among the three baselines based on hand-crafted features, we observe that multi-class classification underperforms both pairwise classification and regression, leaving a negative gap on the small and medium-sized datasets. This highlights that algorithm selection is not a simple multi-class classification problem. In that formulation, the model loses valuable information about the performance of algorithms other than the best one.
Finally, we find more gaps closed for the large dataset across all three baselines, as it includes harder instances. As time increases, the benefit of algorithm selection grows because feature computation accounts for a smaller fraction of the total computation time. Moreover, for harder instances, multiple solvers frequently time out, so choosing a solver that successfully solves the instance becomes especially valuable, given the additional penalty imposed by the \gls{par} score.



\begin{table}[htbp]
	\tabledefaults
	\centering
	\caption[SAT solving algorithm selection results]{\textbf{\gls{sat} Solving.} Algorithm selection results (gap closed; higher is better).}
	\label{table:sat_as_res}
	\resizebox{0.6\textwidth}{!}{%
		\begin{tabular}{llcccc}
			\toprule
			\textbf{Instances} & \textbf{Method} & \dataset{VG}         & \dataset{VCG}        & \dataset{LCG}        & \dataset{SATzilla}    \\
			\midrule
			\multirow{9}{*}{\dataset{Small}}
			                   & GIN             & \meanstd{0.45}{0.00} & \meanstd{0.45}{0.00} & \meanstd{0.45}{0.01}                         \\
			                   & GCN             & \meanstd{0.00}{0.00} & \meanstd{0.40}{0.00} & \meanstd{0.38}{0.03}                         \\
			                   & GCN+            & \meanstd{0.31}{0.00} & \meanstd{0.41}{0.01} & \meanstd{0.47}{0.02}                         \\
			                   & GIN+            & \meanstd{0.28}{0.05} & \meanstd{0.02}{0.01} & \meanstd{0.16}{0.08}                         \\
			                   & Gated GCN+      & --                   & \meanstd{0.40}{0.04} & \meanstd{0.48}{0.03}                         \\
			                   & PR              &                      &                      &                      & \meanstd{0.38}{0.02}  \\
			                   & MC              &                      &                      &                      & \meanstd{-0.44}{0.01} \\
			                   & PC              &                      &                      &                      & \meanstd{-0.26}{0.00} \\
			\midrule
			\multirow{3}{*}{\dataset{Medium}}
			                   & PR              &                      &                      &                      & \meanstd{0.35}{0.01}  \\
			                   & MC              &                      &                      &                      & \meanstd{-0.43}{0.04} \\
			                   & PC              &                      &                      &                      & \meanstd{0.01}{0.01}  \\
			\midrule
			\multirow{3}{*}{\dataset{Large}}
			                   & PR              &                      &                      &                      & \meanstd{0.51}{0.01}  \\
			                   & MC              &                      &                      &                      & \meanstd{-0.05}{0.01} \\
			                   & PC              &                      &                      &                      & \meanstd{0.21}{0.01}  \\
			\bottomrule
		\end{tabular}}
\end{table}
